\section{Problem 15: Erd\H{o}s' alternating prime series $\sum_{n\ge 1} (-1)^n\,\frac{n}{p_n}$ (Round 4)}

\subsection{1) ROUND-4 OBJECTIVE}

\textbf{Path (C): obstruction/correction, strengthened to an exact threshold theorem.}

Round 3 proved a \emph{power-scale} barrier: for discrepancy bounds of the form
\[
A(x):=\sum_{m\le x} \varepsilon(m)=O\!\left(\frac{x}{(\log\log x)^\alpha}\right),
\]
the exponent $\alpha>1$ suffices for convergence of $\sum \varepsilon(m)/(m\log m)$, while every $\alpha\le 1$ is insufficient in general.

The natural next step is to remove the special choice $x/(\log\log x)^\alpha$ and classify \emph{all monotone discrepancy gauges}.  That is what this round does: for any eventually nondecreasing sublinear function $F$, I prove an exact criterion for when a bound
\[
A(x)=O(F(x))
\]
forces convergence of
\[
\sum_{m\ge 3}\frac{\varepsilon(m)}{m\log m}.
\]
This gives a sharp obstruction for every ``discrepancy-only'' strategy, not just the power family from Round 3.

\subsection{2) Round-3 FOUNDATION USED}

I use the following Round-3 results without re-proving them.
\begin{enumerate}
\item \textbf{Round 3, \S5.4 (using Round 2 Theorem 5.7):} the Erd\H{o}s series
\[
\sum_{n\ge 1}(-1)^n\frac{n}{p_n}
\]
converges if and only if the parity series
\[
\sum_{m\ge 2}\frac{(-1)^{\pi(m)}}{m\log m}
\]
converges.
\item \textbf{Round 3, Proposition 5.9 / Corollary 5.3:} the ``sparse odd-set modification of the alternating sequence'' gives divergent examples with controlled discrepancy in the special family $F(x)=x/(\log\log x)^\alpha$.
\end{enumerate}

The new round generalizes that special-family construction to \emph{arbitrary monotone sublinear} gauges $F$.

\subsection{3) NEW INSIGHT / TOOL (ROUND 4)}

\textbf{New tool: an exact threshold theorem for discrepancy-only criteria.}

Let $w_n:=1/(n\log n)$ and let
\[
A_\varepsilon(x):=\sum_{3\le m\le x}\varepsilon(m), \qquad \varepsilon(m)\in\{\pm 1\}.
\]
I will prove:

\begin{quote}
For every eventually nondecreasing function $F:\{3,4,\dots\}\to (0,\infty)$ with $F(n)=o(n)$,
\[
\sum_{n\ge 3}\frac{F(n)}{n^2\log n}<\infty
\]
is \emph{exactly} the threshold for whether the implication
\[
A_\varepsilon(x)=O(F(x)) \implies \sum_{n\ge 3}\frac{\varepsilon(n)}{n\log n}\text{ converges}
\]
holds for all sign sequences $\varepsilon$.
\end{quote}

Thus Round 2's sufficient criterion and Round 3's power-family sharpness are special cases of a single exact theorem.

\subsection{4) ATTACK PLAN (ROUND 4)}

\textbf{Exact gap left by Round 3.}
Round 3 only handled the one-parameter family $F(x)=x/(\log\log x)^\alpha$.  The remaining gap is to determine whether there is a \emph{general} discrepancy-only criterion, and if so, what its exact threshold is.

\textbf{Claims to prove now.}
\begin{enumerate}
\item \textbf{Universal sufficiency:} if $\sum F(n)/(n^2\log n)$ converges and $A_\varepsilon(x)=O(F(x))$, then $\sum \varepsilon(n)/(n\log n)$ converges.
\item \textbf{Universal counterexample:} if $\sum F(n)/(n^2\log n)$ diverges, then there exists an equidistributed sign sequence $\varepsilon$ with $A_\varepsilon(x)=O(F(x))$ for which $\sum \varepsilon(n)/(n\log n)$ diverges.
\end{enumerate}

Together these produce a complete obstruction theorem for discrepancy-only approaches.

\subsection{5) WORK (ROUND 4)}

\subsubsection{5.1 Setup}

Fix a sign sequence $\varepsilon(n)\in\{\pm 1\}$ for $n\ge 3$, and define its discrepancy function
\[
A_\varepsilon(x):=\sum_{3\le m\le x}\varepsilon(m).
\]
Let
\[
w_n:=\frac{1}{n\log n}\qquad (n\ge 3).
\]
Then $w_n$ is positive, decreasing, and $w_n\to 0$.

\subsubsection{5.2 Universal sufficiency by Abel summation}

\paragraph{Lemma 5.1 (Abel criterion for weighted sign series).}
Let $F:\{3,4,\dots\}\to (0,\infty)$ be eventually nondecreasing and satisfy $F(n)=O(n)$.  If
\[
\sum_{n\ge 3}\frac{F(n)}{n^2\log n}<\infty,
\]
and if $A_\varepsilon(x)=O(F(x))$, then the series
\[
\sum_{n\ge 3}\frac{\varepsilon(n)}{n\log n}
\]
converges.


\emph{Proof.}
For integers $3\le M\le N$, Abel summation gives
\[
\sum_{n=M}^{N}\varepsilon(n)w_n
= A_\varepsilon(N)w_N - A_\varepsilon(M-1)w_M + \sum_{n=M}^{N-1} A_\varepsilon(n)\,(w_n-w_{n+1}).
\tag{5.1}
\]

We claim that
\[
0<w_n-w_{n+1}\ll \frac{1}{n^2\log n} \qquad (n\ge 3).
\tag{5.2}
\]
Indeed, with $f(x)=1/(x\log x)$, the mean value theorem gives
\[
w_n-w_{n+1} = -f'(\xi_n)
\]
for some $\xi_n\in(n,n+1)$, and
\[
-f'(x)=\frac{\log x+1}{x^2\log^2 x}
\ll \frac{1}{x^2\log x}
\]
for $x\ge 3$, proving (5.2).

Now assume $|A_\varepsilon(n)|\le C F(n)$ for all large $n$.  Since $F(n)=O(n)$,
\[
|A_\varepsilon(N)w_N| \ll \frac{F(N)}{N\log N}\to 0,
\qquad
|A_\varepsilon(M-1)w_M| \ll \frac{F(M)}{M\log M}\to 0.
\tag{5.3}
\]
Also, by (5.2),
\[
\sum_{n=M}^{\infty}|A_\varepsilon(n)|\,(w_n-w_{n+1})
\ll \sum_{n=M}^{\infty}\frac{F(n)}{n^2\log n}\to 0
\qquad (M\to\infty)
\tag{5.4}
\]
by hypothesis.

Thus the right-hand side of (5.1) tends to $0$ uniformly in $N\ge M$ as $M\to\infty$, so the tails of $\sum \varepsilon(n)w_n$ are Cauchy. Hence the series converges.
\hfill$\square$

\subsubsection{5.3 Dyadic condensation for the threshold series}

\paragraph{Lemma 5.2 (Dyadic condensation).}
Let $F:\{3,4,\dots\}\to (0,\infty)$ be eventually nondecreasing. Then
\[
\sum_{n\ge 3}\frac{F(n)}{n^2\log n}<\infty
\quad\Longleftrightarrow\quad
\sum_{k\ge 2}\frac{F(2^k)}{2^k k}<\infty.
\]
The same equivalence holds with ``$<\infty$'' replaced by ``$=\infty$''.


\emph{Proof.}
For $k\ge 2$, set
\[
B_k:=\sum_{n=2^k}^{2^{k+1}-1}\frac{F(n)}{n^2\log n}.
\]
Because $F$ is nondecreasing, for $2^k\le n<2^{k+1}$ we have
\[
F(2^k)\le F(n)\le F(2^{k+1}).
\]
Also there are exactly $2^k$ integers in the block $[2^k,2^{k+1})$, and throughout that block
\[
2^k\le n<2^{k+1},
\qquad
k\log 2\le \log n\le (k+1)\log 2.
\]
Hence
\[
B_k
\ge
2^k\cdot \frac{F(2^k)}{2^{2k+2}(k+1)\log 2}
= c_1\frac{F(2^k)}{2^k k}
\tag{5.5}
\]
for some absolute constant $c_1>0$, and similarly
\[
B_k
\le
2^k\cdot \frac{F(2^{k+1})}{2^{2k}k\log 2}
= c_2\frac{F(2^{k+1})}{2^k k}
\le c_3\frac{F(2^{k+1})}{2^{k+1}(k+1)}
\tag{5.6}
\]
for absolute constants $c_2,c_3>0$.

Therefore $B_k$ is squeezed between constant multiples of two consecutive dyadic terms.  Summing over $k$ proves that
\[
\sum_{n\ge 3}\frac{F(n)}{n^2\log n}
= \sum_{k\ge 2} B_k
\]
converges if and only if
\[
\sum_{k\ge 2}\frac{F(2^k)}{2^k k}
\]
converges.
\hfill$\square$

\subsubsection{5.4 Universal counterexample when the threshold diverges}

\paragraph{Proposition 5.3 (General sparse-block counterexample).}
Let $F:\{3,4,\dots\}\to (0,\infty)$ be eventually nondecreasing and satisfy $F(n)=o(n)$.  Assume that
\[
\sum_{n\ge 3}\frac{F(n)}{n^2\log n}=\infty.
\]
Then there exists an equidistributed sign sequence $\varepsilon(n)\in\{\pm 1\}$ such that
\[
A_\varepsilon(x)=O(F(x))
\]
and yet
\[
\sum_{n\ge 3}\frac{\varepsilon(n)}{n\log n}
\]
diverges.


\emph{Proof.}
By Lemma 5.2,
\[
\sum_{k\ge 2}\frac{F(2^k)}{2^k k}=\infty.
\tag{5.7}
\]
Fix a constant $c\in(0,1/10)$ and define
\[
G_k:=\bigl\lfloor cF(2^k)\bigr\rfloor \qquad (k\ge 2).
\]
Since $F$ is eventually nondecreasing, the sequence $(G_k)$ is eventually nondecreasing.  Also, because $F(n)=o(n)$,
\[
G_k=o(2^k).
\tag{5.8}
\]
After discarding finitely many initial $k$, we may assume for every $k\ge k_0$ that
\[
G_k\le 2^{k-2}.
\tag{5.9}
\]

Define the increments
\[
M_{k_0}:=G_{k_0},
\qquad
M_k:=G_k-G_{k-1} \quad (k>k_0).
\]
Then $M_k\ge 0$ for all $k\ge k_0$.

For each $k\ge k_0$, let
\[
I_k:=[2^k,2^{k+1})\cap \mathbb{N}.
\]
Choose exactly $M_k$ odd integers from $I_k$; this is possible by (5.9), because $I_k$ contains $2^{k-1}$ odd integers.  Let $Y$ be the union of all chosen odd integers.

Now define
\[
\varepsilon(n):=
\begin{cases}
+1, & n\in Y,\\
(-1)^n, & n\notin Y.
\end{cases}
\tag{5.10}
\]
So we start from the alternating sequence $(-1)^n$ and flip the sign at selected odd positions.

\paragraph{Step 1: discrepancy bound.}
Let $N\ge 2^{k_0}$ and let $K$ be such that $2^K\le N<2^{K+1}$.  The baseline sequence $(-1)^n$ has discrepancy bounded by $1$, and each modified odd index changes the partial sum by exactly $+2$.  Therefore
\[
|A_\varepsilon(N)|
\le 1 + 2\#(Y\cap[1,N]).
\tag{5.11}
\]
But by construction,
\[
\#(Y\cap[1,N])\le \sum_{k_0\le j<K}M_j + M_K = G_K.
\tag{5.12}
\]
Since $F$ is nondecreasing and $N\ge 2^K$,
\[
G_K\le cF(2^K)\le cF(N).
\tag{5.13}
\]
Hence
\[
A_\varepsilon(N)=O(F(N)).
\tag{5.14}
\]
Because $F(N)=o(N)$, we also get
\[
A_\varepsilon(N)=o(N),
\tag{5.15}
\]
so the sequence $\varepsilon$ is equidistributed.

\paragraph{Step 2: divergence of the weighted series.}
We decompose
\[
\sum_{n\ge 3}\frac{\varepsilon(n)}{n\log n}
=
\sum_{n\ge 3}\frac{(-1)^n}{n\log n}
+
\sum_{n\in Y}\frac{\varepsilon(n)-(-1)^n}{n\log n}.
\tag{5.16}
\]
The first series converges by the alternating series test.  For every $n\in Y$ we have $n$ odd, hence $(-1)^n=-1$ and $\varepsilon(n)=+1$, so
\[
\varepsilon(n)-(-1)^n = 2.
\tag{5.17}
\]
Thus the second series equals
\[
2\sum_{n\in Y}\frac{1}{n\log n}.
\tag{5.18}
\]
It remains to show that this diverges.

For $n\in I_k$ one has $n<2^{k+1}$ and $\log n\le (k+1)\log 2$, so every chosen $n\in Y\cap I_k$ satisfies
\[
\frac{1}{n\log n}
\ge \frac{1}{2^{k+1}(k+1)\log 2}
\gg \frac{1}{2^k k}.
\tag{5.19}
\]
Therefore
\[
\sum_{n\in Y}\frac{1}{n\log n}
\gg \sum_{k\ge k_0}\frac{M_k}{2^k k}.
\tag{5.20}
\]
Set
\[
q_k:=\frac{1}{2^k k}.
\]
Then summation by parts gives, for $K\ge k_0$,
\[
\sum_{k=k_0}^{K} M_k q_k
=
G_K q_K - G_{k_0-1}q_{k_0} + \sum_{k=k_0}^{K-1} G_k(q_k-q_{k+1}).
\tag{5.21}
\]
Now
\[
q_k-q_{k+1}
= \frac{1}{2^k k}-\frac{1}{2^{k+1}(k+1)}
= \frac{k+2}{2^{k+1}k(k+1)}
\ge \frac12\,q_k,
\tag{5.22}
\]
and by (5.8),
\[
G_K q_K \le \frac{cF(2^K)}{2^K K}\to 0.
\tag{5.23}
\]
So if
\[
\sum_{k\ge k_0} G_k q_k = \infty,
\tag{5.24}
\]
then (5.21)--(5.23) imply
\[
\sum_{k\ge k_0}\frac{M_k}{2^k k}=\infty.
\tag{5.25}
\]
Finally,
\[
G_k \ge cF(2^k)-1,
\tag{5.26}
\]
and since $\sum 1/(2^k k)<\infty$, equation (5.7) implies
\[
\sum_{k\ge k_0}\frac{G_k}{2^k k}=\infty.
\tag{5.27}
\]
Hence (5.24) holds, so by (5.25), (5.20), and (5.18), the series $\sum \varepsilon(n)/(n\log n)$ diverges.
\hfill$\square$

\subsubsection{5.5 Exact threshold theorem}

\paragraph{Theorem 5.4 (Exact threshold for discrepancy-only methods).}
Let $F:\{3,4,\dots\}\to (0,\infty)$ be eventually nondecreasing and satisfy $F(n)=o(n)$.  Then the following are equivalent:
\begin{enumerate}
\item \label{it:threshold-sum}
\[
\sum_{n\ge 3}\frac{F(n)}{n^2\log n}<\infty.
\]
\item \label{it:all-signs}
For every sign sequence $\varepsilon(n)\in\{\pm 1\}$ with $A_\varepsilon(x)=O(F(x))$, the weighted series
\[
\sum_{n\ge 3}\frac{\varepsilon(n)}{n\log n}
\]
converges.
\end{enumerate}
Moreover, if \ref{it:threshold-sum} fails, then there exists an equidistributed sign sequence $\varepsilon(n)$ with $A_\varepsilon(x)=O(F(x))$ for which the weighted series diverges.


\emph{Proof.}
The implication \ref{it:threshold-sum}$\Rightarrow$\ref{it:all-signs} is Lemma 5.1.  The converse obstruction is Proposition 5.3.
\hfill$\square$

\subsubsection{5.6 Return to Erd\H{o}s problem 15}

Define the prime-parity discrepancy
\[
A_\pi(x):=\sum_{3\le m\le x}(-1)^{\pi(m)}.
\]
Round 3 (via Round 2 Theorem 5.7) established that the Erd\H{o}s series converges if and only if
\[
\sum_{m\ge 3}\frac{(-1)^{\pi(m)}}{m\log m}
\]
converges.  Combining that fixed Round-3 foundation with Theorem 5.4 yields the following corrected formulation.

\paragraph{Corollary 5.5 (Minimal corrected discrepancy criterion).}
Let $F$ be eventually nondecreasing with $F(n)=o(n)$.
\begin{enumerate}
\item If
\[
A_\pi(x)=O(F(x))
\quad\text{and}\quad
\sum_{n\ge 3}\frac{F(n)}{n^2\log n}<\infty,
\]
then the Erd\H{o}s series
\[
\sum_{n\ge 1}(-1)^n\frac{n}{p_n}
\]
converges.
\item If
\[
\sum_{n\ge 3}\frac{F(n)}{n^2\log n}=\infty,
\]
then a proof of convergence based \emph{only} on a discrepancy bound of the form $A_\pi(x)=O(F(x))$ cannot be valid in general: there exist equidistributed sign sequences with the same size bound whose weighted series diverges.
\end{enumerate}


This is the exact barrier statement sought in this round.

\subsubsection{5.7 Concrete refinement beyond Round 3}

Round 3 treated only the family $F(x)=x/(\log\log x)^\alpha$.  Theorem 5.4 gives many sharper examples immediately.

\paragraph{Corollary 5.6 (A stricter log--log--log threshold).}
Fix $r\in\mathbb{R}$ and, for $x$ large, define
\[
F_r(x):=\frac{x}{(\log\log x)(\log\log\log x)^r}.
\]
Then:
\begin{enumerate}
\item If $r>1$, every sign sequence $\varepsilon(n)$ with
\[
A_\varepsilon(x)=O(F_r(x))
\]
has convergent weighted series $\sum \varepsilon(n)/(n\log n)$.
\item If $r\le 1$, there exists an equidistributed sign sequence $\varepsilon(n)$ with
\[
A_\varepsilon(x)=O(F_r(x))
\]
for which $\sum \varepsilon(n)/(n\log n)$ diverges.
\end{enumerate}


\emph{Proof.}
Apply Theorem 5.4 with $F=F_r$.  The threshold series is
\[
\sum_{n\ge 3}\frac{F_r(n)}{n^2\log n}
\asymp
\sum_{n\ge 3}\frac{1}{n\log n\log\log n\,(\log\log\log n)^r},
\]
which converges if and only if $r>1$ by the integral test with the substitutions $u=\log\log n$ and $v=\log u$.
\hfill$\square$

This strict refinement was not visible in Round 3, so it is genuine new progress.

\subsection{6) ADVERSARIAL VERIFICATION}

\begin{itemize}
\item \textbf{Did I accidentally assume more regularity than stated?} No.  The threshold theorem uses only: (i) eventual monotonicity of $F$, for dyadic condensation and counting control; and (ii) $F(n)=o(n)$, for equidistribution and to guarantee enough room in each dyadic block.
\item \textbf{Is Lemma 5.2 valid without a doubling hypothesis?} Yes.  The upper bound uses the shifted endpoint value $F(2^{k+1})$, but this only shifts the dyadic series by one index and a bounded factor; no doubling condition is required.
\item \textbf{Could the floor in $G_k=\lfloor cF(2^k)\rfloor$ destroy divergence?} No.  Since $G_k\ge cF(2^k)-1$ and $\sum 1/(2^k k)<\infty$, subtracting the floor error cannot change divergence of $\sum G_k/(2^k k)$.
\item \textbf{Are there enough odd numbers in each block?} Yes.  Since $F(2^k)=o(2^k)$, eventually $G_k\le 2^{k-2}$, while the block $[2^k,2^{k+1})$ contains $2^{k-1}$ odd numbers.
\item \textbf{Does the theorem say anything directly about the prime-parity sequence?} Not by itself.  It classifies what can be concluded from a \emph{discrepancy-only} estimate $A_\pi(x)=O(F(x))$.  The actual prime sequence may have extra structure; the theorem rules out an entire class of proofs, but does not solve the Erd\H{o}s problem unconditionally.
\end{itemize}

\subsection{7) FINAL (EXACTLY ONE)}

\textbf{UNRESOLVED (BUT STRICTLY ADVANCED).}

This round proves a full abstract theorem that strictly strengthens Round 3:
\begin{itemize}
\item Round 2 gave a sufficient discrepancy criterion.
\item Round 3 showed sharpness for the special family $F(x)=x/(\log\log x)^\alpha$.
\item Round 4 gives the \emph{exact threshold for every monotone sublinear gauge $F$}: the discrepancy-only implication
\[
A(x)=O(F(x))\Longrightarrow \sum \frac{\varepsilon(n)}{n\log n}\text{ converges}
\]
holds for all sign sequences if and only if
\[
\sum \frac{F(n)}{n^2\log n}<\infty.
\]
\end{itemize}

For Erd\H{o}s problem 15, this isolates the remaining issue sharply: any unconditional proof based only on a bound for
\[
A_\pi(x)=\sum_{m\le x}(-1)^{\pi(m)}
\]
must beat the exact threshold above; anything weaker is provably insufficient in general.

\subsection{8) COMPLETION ESTIMATE (MANDATORY)}

\noindent\textbf{COMPLETION: 82\%}

\subsection{9) REFERENCES}

\begin{thebibliography}{9}
\bibitem{EP15}
Erd\H{o}s Problems website,
\emph{Problem 15: Is it true that $\sum_{n=1}^\infty (-1)^n n/p_n$ converges?}
\texttt{https://www.erdosproblems.com/latex/15}.

\bibitem{Tao2024}
T. Tao,
\emph{The convergence of an alternating series of Erd\H{o}s, assuming the Hardy--Littlewood prime tuples conjecture},
Communications of the American Mathematical Society \textbf{4} (2024), 80--96.

\bibitem{TaoArxiv}
T. Tao,
\emph{The convergence of an alternating series of Erd\H{o}s, assuming the Hardy--Littlewood prime tuples conjecture},
arXiv:2308.07205.
\end{thebibliography}

\subsection{10) OUTPUT FORMAT}

The Erd\H{o}s problem number here is \textbf{15}.  This file is therefore exported as \texttt{15\_solutions.tex}, beginning directly with \verb|\section| and with no preamble.
